\section{The selected bridge and gauge-fixing map}
\label{sec:ex-selected-stationary-linear-data}

The preceding estimates must be supplied on the actual retained horizon.
We verify two identifications: the fine-difference bounds of its linear
stationary bridge, and the exact gauge identity of its fluctuation matrix.

\subsection{The unchanged one-form transport}
Let $A_0$ be the one-form transport of the retained tower at a fixed
adjacent pair of scales. The construction in
Section~\ref{sec:curvature-horizon-producer} retains this row unchanged;
only the curvature completion rows are replaced. Apply
Theorem~\ref{thm:ex-lattice-wire-tower-strip} with $e=1$. The schedule
has $j\geq4$, so every difference word $u$ of length at most two is
covered. Its actual Laurent row satisfies
\[
       \sum_z e^{w|z|_1}|(\nabla^u A_0)(z)|
                         \leq D L^{-1-|u|}.
\]
The width $w>0$ and the cost $D$ are those of the stage of this
adjacent pair. They are common to every terminal cutoff, every level
at that stage and every finite period, but, as stated at the end of
Section~\ref{sec:curvature-horizon-producer}, no common positive width
or common cost over all stages is asserted. The hypotheses of the cited
theorem for the retained rows, and the resulting difference bounds, are
part of the horizon supplier of that section; here they are an input.
The price enlargement used in the retained construction can only
increase $D$ and narrow $w$.

Define $\Phi(p)$ by Fourier evaluation of $A_0$, with the color identity
on its color factor. Fourier evaluation commutes with a cover difference,
including its centered carry: this is exactly the residue-difference
identity in the proof of the cited theorem. Thus, for a coordinate leg,
\[
       |\nabla^u\Phi(p)e_a(z)|\leq D L^{-1-|u|}
                      \qquad (p\in\operatorname{strip}(w),\ |u|\leq2).
\]
Absolute exponential summability gives holomorphy and periodicity on the
strip. Restriction to a smaller width preserves every bound. This proves
\eqref{eq:ex-vertex-cruise} for the selected bridge and supplies the
identical hypothesis in the inherited-loop estimate.

We must also identify this row with the derivative of the stationary
representative used below. Write $W=A_0$ in finite-volume coordinates.
In the original tree gauge chart the positive internal Hessian gives a
unique local stationary branch. Its fine tangent is not $W$ itself:
straightening up the tree adds an infinitesimal gauge transformation,
so the expected tangent is
\[
                         E=W+d_0K,
\]
where $K$ is the linear tree-straightening map: at a fine site, $K$
integrates $W$ along the tree path to the block root and subtracts the
root value $\operatorname{fixRoot}W$. Existence and uniqueness of the
branch do not by themselves show that its derivative is this particular
matrix. The identification is the following input, which is the
installed-tangent identity proved in
Section~\ref{subsec:path-stationary-chart}.

\begin{quote}
\emph{Stationary tangent input.} On each finite torus of the selected
family, for the same forest, incoming path, fine and coarse beds and the
same installed matrix $W$ of the retained tower, the stationary branch
$\rho$ of the tree chart satisfies
\[
          D\rho(0)x=(Ex-\iota x)\big|_{\rm free}
\]
for every coarse direction $x$. Here the tree chart records the fine
field on the free (non-tree) links relative to the linear copy
$\iota x$ of the coarse field, which equals the coarse value on fine
links joining two different blocks and zero on links inside a block.
\end{quote}

\noindent
This is the precise form of the statement that the fine tangent of
the branch is $E$. The cited section derives it by differentiating the
vertical stationarity equation: by the finite-volume Hessian identity
the bilinear form in that equation is the fine bed kinetic form plus
the penalty pairing of the predictor residual, the candidate $E$ solves
the same linear equation by the two matching identities
$K_fW=Q^{\mathsf T}K_c$ and $QW=I-\kappa^{-1}K_c$ of the retained tower
for these two beds, and positivity of the vertical Hessian gives
uniqueness. Its premises (the Hessian identity, the Hessian gap and the
two matching identities for the same beds and the same $W$) are
therefore premises of everything that follows. Gauge this branch back
by the group exponential of $Kx$. The derivative of the resulting fine
field is $E-d_0K=W$. At a block root the straightening identity reads
$K_{\rm root}=-\operatorname{fixRoot}W$. The derivative of the rooted
frame is $-\operatorname{fixRoot}W$ as well; multiplication by the inverse
root gauge factor cancels it. Thus recentering the gauged-back branch
gives the joint derivative $(W,0)$. Gauge covariance preserves its
stationarity. We use uniqueness only in the original tree chart:
the centered merit on the full joint space has gauge null directions.
Finite Fourier inversion now identifies this actual derivative on every
torus grid with $\Phi$. Hence the difference bounds above concern the
selected stationary representative.

\subsection{The exact gauge identity}
Work first on one finite torus. Let $d$ be the infinitesimal fine gauge
action at the identity, in the sign convention of the joint action,
and let $\mathcal R$ be its expression in centered block coordinates.
Its two components are $(-d_0a,-Ja)$ as in the preceding section.
Let $X$ be the negative right logarithmic derivative of the frame in
these coordinates, and let $C_\mu$ be the positive regulator with scalar
symbol $c_\mu$. The linear gauge-fixing map is
\[
                         F=d^*\operatorname{pr}_{\rm fine}+C_\mu X.
\]
For a gauge curve through the identity, frame covariance gives its frame
ratio as the inverse gauge transformation. Applying the negative right
logarithm therefore returns exactly the gauge parameter near zero.
Differentiating this identity gives $X\mathcal R=I$, after expressing
the gauge tangent in the same coordinates. Also the fine component of
the gauge tangent is $d$. Consequently
\begin{equation}\label{eq:ex-actual-gauge-identity}
                           F\mathcal R=d^*d+C_\mu=\Gamma_\mu.
\end{equation}

There is a coordinate issue in this calculation. The rooted gauge
coordinates and the centered coordinates are related by a triangular
change whose derivative is
\[
                  (u,r)\longmapsto(u,r+K_0u),
\]
where $K_0$ is the derivative of the recentering frame. The frame
covariance identity gives $K_0d_0a=a_{\rm root}-Ja$.
It follows by substitution that this triangular derivative sends the
centered gauge tangent to the rooted gauge tangent. The fine projection
is unchanged. The chain rule sends $X$ to its product with that same
triangular derivative, so both sides of
\eqref{eq:ex-actual-gauge-identity} remain identical. This verifies the
identity in the centered chart used for the estimates, including its
sign; it would be incorrect to combine the rooted $X$ with a centered
$\mathcal R$ without this change.

Fourier transformation yields the identity on every torus grid. The
symbols are common to the unbounded family of periods, by the local
stencil construction. They therefore agree at every real momentum by
density and at every complex momentum in their common strip by analytic
uniqueness. For real momentum, the Fourier symbol of each real finite-volume
matrix satisfies $F(-p)=\overline{F(p)}$. This applies to the frame
linearization, the fine difference, and the real regulator, hence to
$F$ itself. Density of the common finite Fourier grids gives this
identity at every real momentum, and therefore $F(-p)^T=F(p)^*$ there.
Thus the bilinear holomorphic continuation used below agrees with the
positive Hermitian square on the real slice.
The same argument identifies the constant stationary matrix
with $\mathcal G+F(-p)^TF(p)$: the fine action contributes its attached
kinetic form, the block penalty contributes $\kappa N(-p)^TN(p)$, and
the gauge-fixing term contributes the displayed square. Each is the
Hessian of the corresponding term of the same action at its identity.

\subsection{The uniform bound for the gauge-fixing map}
For a unit vector in $W_p$, the fine derivative term $d^*u$ is bounded
in unnormalized Euclidean norm by a constant depending only on $G$,
by the Fourier multiplier estimate in
Section~\ref{sec:ex-weighted-frame-penalty}. The frame derivative $Xv$
has normalized block $L^2$ norm at most a constant, by the centered-frame
induction, restriction to one residue block, and the finite coordinate
operations defining its logarithm. Restriction from the padded patch
to this block costs a bounded ratio of volumes. Its unnormalized
Euclidean norm is therefore at most $C L^2$.
On the fixed small coarse strip the regulator has Euclidean operator
norm at most $C'\mu^2$: its alias symbol is
$\mu^2(\mu^2/(\omega+\mu^2))^9$, with the denominator estimate of the
preceding section, and the two Fourier conjugations have uniformly
bounded norms. Hence
\[
                       |C_\mu Xv|\leq C'C\mu^2L^2=C'C.
\]
Adding the fine derivative term gives the required bound $|Fv|\leq D$,
with $D$ fixed before $L$. The same stencil construction is holomorphic
and periodic in the external momentum, since it uses only linear Fourier
legs in finite real derivatives at the identity and the holomorphic
regulator symbol. This proves both parts of
\eqref{eq:ex-gauge-fixing-receipt} for the selected data.

The inverse estimates of
Section~\ref{sec:ex-weighted-stationary-inverses} now apply to the actual
common stationary matrix. At the origin the two stationary matrix
presentations used for the action and its comparison have this same
matrix. Their inverse bounds consequently use the same width and cost;
no bound for their positive-order derivatives is asserted by this
constant-matrix identification.
